原则上，Kirchhoff 定律足以处理绝大多数集中参数电路；但网络一旦变复杂，直接设置每条支路电流、
列出全部节点和回路方程，会迅速产生规模很大的线性方程组。这样的做法可靠，却冗长、难以检查，也很
容易在符号上出错。因此，对只含线性元件的网络，应优先利用线性叠加、等效电源、对称性、电桥平衡、
星形--三角形变换和自相似性，把网络先化简，再在必要时使用 Kirchhoff 方程。对于计算机，求解大规模
线性方程组并不困难；对人工竞赛解题而言，找到结构才是关键。

\subsection*{叠加原理}

节点电压法、支路电流法和 Kirchhoff 方程不仅适用于电阻，也可与任意元件的伏安关系联立。若元件是
非线性的，得到的通常是非线性方程组，叠加原理一般不成立；若网络只含线性电阻和独立电源，方程组对
电源产生的非齐次项具有可加性。于是多个独立电源同时工作时，任一支路电流或任一节点电势，等于各个
电源单独作用时对应结果的代数和。

考察某一个独立电源的贡献时，其余独立电源要“置零”，同时保留它们的内阻和原有连接关系：

\begin{itemize}
  \item 理想电压源置零，即把其电压改为零，等效为短路；
  \item 理想电流源置零，即把其电流改为零，等效为断路。
\end{itemize}

下图把上一讲含电压源 $E_0$、电流源 $I_0$ 和两个电阻的电路分成两个子问题。仅保留电压源时，
$R_1,R_2$ 串联，流过 $R_2$ 的电流为
\[
  I_1=\frac{E_0}{R_1+R_2}.
\]
仅保留电流源时，$R_1,R_2$ 并联，按分流关系，流过 $R_2$ 的电流为
\[
  I_2=\frac{R_1}{R_1+R_2}I_0.
\]
若图示两部分在 $R_2$ 中产生的电流方向相同，则原电路中
$I_{R_2}=I_1+I_2$；若方向相反，就应按所选正方向作代数相加。

\superpositionFigure
\diagramcaption{图 6.1：叠加原理。左图只保留电压源，右图只保留电流源；被置零的电压源短路，被置零的电流源断路。}

叠加的是电压和电流，不是功率。由于 $P=I^2R$ 或 $P=U^2/R$ 是二次量，通常不能把各电源单独作用时
的功率直接相加。

\subsection*{等效电源定理}

线性叠加还有两个重要推论：戴维南（Thevenin）等效电压源定理和诺顿（Norton）等效电流源定理。
无论内部结构多复杂，一个线性有源二端网络对外只表现为端电压与端电流之间的一条直线关系；只要确定
这条直线的截距和斜率，就不必再追踪网络内部的每条支路。

对图 6.2(a) 的二端网络，规定测试电流 $i_t$ 流入上端口，端电压为 $v_t$。在端口接入可调理想电流源，
得到图 6.2(b)。由叠加原理，$v_t$ 可拆成两部分：

\begin{enumerate}
  \item 保留测试电流源，把网络内部所有独立源置零。此时内部成为纯电阻网络，端口等效电阻记为 $r$，
  因而图 6.2(c) 中 $v_a=i_t r$；
  \item 断开测试电流源，恢复网络内部全部独立源。端口电流为零，图 6.2(d) 中的端电压 $v_b$ 就是
  开路电压。
\end{enumerate}

所以
\[
  v_t=v_a+v_b=i_t r+v_b.
\]
这正是一个理想电压源与内阻串联时的端口关系。

\theveninDerivationFigure
\diagramcaption{图 6.2：戴维南定理的推导。四幅图的端口方向完全一致，第三幅只保留测试源，第四幅只保留网络内部独立源。}

具体地说，任意线性有源二端网络都可等效为开路电压 $V_{\rm Th}$ 与电阻 $R_{\rm Th}$ 串联：
$V_{\rm Th}$ 由原网络端口开路时的电压得到；$R_{\rm Th}$ 由内部所有独立源置零后、从端口向里看见的
等效电阻得到。若网络中还含受控源，受控源不能置零，应在端口加测试源求 $R_{\rm Th}=v_t/i_t$。

同一网络也可等效为短路电流 $I_{\rm N}$ 与电阻 $R_{\rm N}$ 并联。$I_{\rm N}$ 是端口短路时的电流，
$R_{\rm N}$ 与 $R_{\rm Th}$ 相同。两种等效电路在端口上满足
\begin{equation}
  \boxed{R_{\rm Th}=R_{\rm N},\qquad V_{\rm Th}=I_{\rm N}R_{\rm N}.}
\end{equation}

\sourceEquivalenceFigure
\diagramcaption{图 6.3：同一个线性有源二端网络的戴维南形式与诺顿形式。等效只保证端口外部特性相同。}

“等效”并不意味着三种电路内部的电流和功率分布相同；它只表示接到同一个负载后，端口电压和负载电流
完全一致。


% EM teaching enrichment 20260927
\paragraph{负载功率为何在阻值相等时最大}
当等效内阻 $R_{\rm Th}>0$、外接纯电阻 $R_L>0$ 时，
\[
 P_L=\frac{V_{\rm Th}^2R_L}{(R_{\rm Th}+R_L)^2},\qquad
 \frac{\dd P_L}{\dd R_L}
 =\frac{V_{\rm Th}^2(R_{\rm Th}-R_L)}{(R_{\rm Th}+R_L)^3}.
\]
故 $R_L=R_{\rm Th}$ 时最大，$P_{L,\max}=V_{\rm Th}^2/(4R_{\rm Th})$。这里最大化的是负载功率，不是输电效率；在这个简单串联等效模型中，匹配时负载与内阻耗散相等，效率仅为一半。端口等效也不能据此还原复杂原网络内部每个元件的耗散。

\subsection*{对称性}

若交换网络中的两个节点后，所有连接、元件参数以及外加电源条件都保持不变，那么这两个节点不可区分，
电势必然相等。等势节点之间没有电流，可以按需要用理想导线相连，也可以把连接它们的支路删去。使用
对称性时，必须把“电路本身的对称”与“端口激励的对称”同时检查；只看几何外形相似而忽略电流从哪里
流入、从哪里流出，是复杂电阻网络中最常见的错误之一。


% EM teaching enrichment 20260927
\paragraph{等势结论来自同一个边值问题}
把端口 $A$ 固定为 $U$、$B$ 固定为零。若一个变换保持两个端口各自不动，并把等阻支路一一对应，变换前后的节点电势都满足同一组 KCL；由电阻网络解的唯一性，相互交换的节点必等势。
若变换交换 $A,B$，一般得到的是 $V_{i'}=U-V_i$，而非 $V_{i'}=V_i$。只有处在该变换不动位置的节点才能推出 $V=U/2$。
例如所有棱阻值为 $R$ 的正方体，从一对体对角顶点通电，靠近入端的三个顶点等势，靠近出端的三个也等势，三层等效为
\[
 R_{AB}=\frac R3+\frac R6+\frac R3=\frac{5R}{6}.
\]
若把出端换到相邻顶点，原先的三点分组不再保持端口条件，不能照搬上述结果。

\subsection*{电桥平衡}

图 6.4 中，电源跨接在左右两端。若桥臂满足
\[
  \boxed{R_1R_4=R_2R_3},
\]
则上、下两个中点电势相等，$R_5$ 中电流为零，这就是电桥平衡。此时 $R_5$ 既可看作没有电流的断路支路
而删去，也可利用两端等势把它短接；两种处理对其余电路的电流、电压都没有影响。要注意，“断路”说的是
$R_5$ 中无电流，“短路”说的是它的两个端点等势，二者描述的是同一平衡状态的不同侧面。

\bridgeFigure
\diagramcaption{图 6.4：平衡电桥。$R_5$ 的上下端分别是两个分压支路的中点；平衡条件下二者等势。}


% EM teaching enrichment 20260927
\paragraph{先断开桥支路，比较两个分压}
设图中上支路由左至右为 $R_1,R_2$，下支路为 $R_3,R_4$，左右端电势分别为 $U,0$。暂断开 $R_5$，两中点电势为
\[
 V_{\rm up}=U\frac{R_2}{R_1+R_2},\qquad
 V_{\rm down}=U\frac{R_4}{R_3+R_4}.
\]
相等即给出 $R_1R_4=R_2R_3$。接回任意有限 $R_5$ 不产生桥电流，因而原分压保持成立；若不相等，接回后两个分压都会受加载影响。
还可把桥端口外的网络化成戴维南源，其内阻为
$(R_1\parallel R_2)+(R_3\parallel R_4)$，桥电流便等于开路中点电压差除以该内阻与 $R_5$ 之和。这给出了偏离平衡时可继续计算的路线。

\subsection*{电流分布法}

求两点间等效电阻时，常规做法是在两点间加电压 $U$，求出总电流 $I$，再用 $R_{\rm eq}=U/I$。电流
分布法则反过来：从一点灌入给定电流 $I$，从另一点抽出同样的电流，先利用对称性和 KCL 确定各支路
如何分流，再求两点间电势差 $U$，最后用 $R_{\rm eq}=U/I$。

这种方法同样服从叠加。例如，要得到“电流从 $A$ 进、从 $B$ 出”的分布，可以把“从 $A$ 进、从
$C$ 出”和“从 $C$ 进、从 $B$ 出”两种分布按相同电流强度相加。中间节点 $C$ 的注入与抽出相互抵消，
剩下的正是所求端口条件。


% EM teaching enrichment 20260927
\paragraph{分流必须同时满足回路电压条件}
只写 KCL 可以构造许多看似守恒的支路电流，却未必对应真实电势。得到候选分流后，还应检查任意独立回路的 $\sum IR=0$。
一个实用核验是功率：若从 $A$ 注入 $I$、从 $B$ 抽出 $I$，则
\[
 I(V_A-V_B)=\sum_{\text{各电阻支路}}R_kI_k^2,
 \qquad R_{AB}=\frac{\sum_kR_kI_k^2}{I^2}.
\]
每条支路只计一次。对三角形三个等阻 $R$、端口取任意两顶点，直连支路分得 $2I/3$，绕行两阻各为 $I/3$，代入得 $R_{AB}=2R/3$，与 $R\parallel2R$ 相符。

\subsection*{多端纯电阻网络}

只有 $N$ 个端口能与外界连接、内部只含电阻的电路，称为 $N$ 端纯电阻网络。一般的无源线性 $N$ 端
电阻网络，可以用每两端口之间各接一个适当电阻的完全图表示。端口若具有对称性，若干等效电阻可能相等；
若无对称性，完全图中的各支路阻值一般不同。

三端网络是最重要的特例。它既可写成三角形（$\Delta$）连接，也可写成星形（$Y$）连接。设端口电势
为 $\varphi_a,\varphi_b,\varphi_c$，星形中心电势为 $\varphi_0$，三条星形支路为 $R_a,R_b,R_c$。
中心节点满足 KCL：
\begin{equation}
  \frac{\varphi_a-\varphi_0}{R_a}
  +\frac{\varphi_b-\varphi_0}{R_b}
  +\frac{\varphi_c-\varphi_0}{R_c}=0.
\end{equation}
解得
\begin{equation}
  \varphi_0=
  \frac{\dfrac{\varphi_a}{R_a}+\dfrac{\varphi_b}{R_b}+\dfrac{\varphi_c}{R_c}}
       {\dfrac1{R_a}+\dfrac1{R_b}+\dfrac1{R_c}}
  =\frac{\varphi_aR_bR_c+\varphi_bR_cR_a+\varphi_cR_aR_b}
  {R_aR_b+R_bR_c+R_cR_a}.
\end{equation}
由端口 $A$ 流入星形网络的电流为
\begin{equation}
  I_a=\frac{\varphi_a-\varphi_0}{R_a}
  =\frac{\varphi_a(R_b+R_c)-\varphi_bR_c-\varphi_cR_b}
  {R_aR_b+R_bR_c+R_cR_a}
  =\frac{\varphi_a-\varphi_b}{R_{AB}}
  +\frac{\varphi_a-\varphi_c}{R_{AC}}.
\end{equation}
最后一个等号是三角形连接的端口电流。令
\[
  S=R_aR_b+R_bR_c+R_cR_a,
\]
比较 $\varphi_a,\varphi_b,\varphi_c$ 的系数，并循环置换端口，可得星形变三角形：
\[
  \boxed{R_{AB}=\frac{S}{R_c},\qquad
  R_{AC}=\frac{S}{R_b},\qquad
  R_{BC}=\frac{S}{R_a}}.
\]
反解得到三角形变星形：
\[
\boxed{
  R_a=\frac{R_{AB}R_{AC}}{R_{AB}+R_{AC}+R_{BC}},\quad
  R_b=\frac{R_{AB}R_{BC}}{R_{AB}+R_{AC}+R_{BC}},\quad
  R_c=\frac{R_{AC}R_{BC}}{R_{AB}+R_{AC}+R_{BC}}.}
\]

\deltaStarFigure
\diagramcaption{图 6.5：同一三端网络的三角形连接与星形连接。对应关系由端口字母决定，不应靠电阻在图中的朝向记忆。}

星形--三角形变换能把许多表面上既不串联也不并联的网络改写成可逐级化简的形式；但在高度对称或
自相似网络中，先用等势、分流和递归往往更简洁。

% EM teaching enrichment 20260927
\paragraph{用三个端口状态检验星三角变换}
若三角形三边均为 $R$，对应星形每臂为 $R/3$。只在 $A,B$ 之间通电、$C$ 外端开路时，星形等效为 $2R/3$，恰与三角形的 $R\parallel2R$ 相同。
但若把 $C$ 接地，它就不再是开路条件，不能把通向 $C$ 的星形臂删掉。三端等效保证任意端电势下的三端电流均相同，比只匹配某一对端口的一个电阻值更强。
复杂网络宜先识别可局部替换的三端块，替换时保留端口名字；随后重新判断串、并联关系，避免按图形朝向误配公式分母。

